\documentclass[11pt,a4paper]{article}

\usepackage[margin=2.5cm]{geometry}
\usepackage{amsmath,amssymb,booktabs,array}
\usepackage{setspace}
\usepackage{natbib}
\usepackage{hyperref}
\begin{document}

\subsection{Translation Invariance and the Cents--Dollars Comparison}

A central finite-sample issue in the application to the Vietnam draft lottery is
whether the estimator is invariant to arbitrary changes in the measurement unit
of the outcome. In the Angrist (1990) application, the outcome is log wages.
Let \(w_i\) denote wages measured in dollars. If wages are instead measured in
cents, the corresponding wage variable is \(100w_i\). Therefore, the log outcome
changes as follows:
\[
    \log(100w_i)
    =
    \log(w_i) + \log(100).
\]
Thus, switching from log wages in dollars to log wages in cents does not change
the economic outcome. It only adds the same constant,
\[
    c = \log(100),
\]
to every observation:
\[
    Y_i^{cents} = Y_i^{dollars} + c.
\]

This is important because a causal effect estimator should not depend on the
arbitrary unit in which wages are recorded before taking logs. If an estimator
changes when wages are measured in cents rather than dollars, the estimate
reflects not only the causal design but also an arbitrary measurement convention.

To see this formally, consider an estimator that can be written as a weighted
sum of outcomes:
\[
    \hat{\tau}
    =
    \sum_{i=1}^n \omega_i Y_i.
\]
If the outcome is shifted by a constant \(c\), then
\[
    \hat{\tau}(Y+c)
    =
    \sum_{i=1}^n \omega_i (Y_i+c).
\]
Expanding the expression gives
\[
    \hat{\tau}(Y+c)
    =
    \sum_{i=1}^n \omega_i Y_i
    +
    c \sum_{i=1}^n \omega_i.
\]
Hence,
\[
    \hat{\tau}(Y+c)
    -
    \hat{\tau}(Y)
    =
    c \sum_{i=1}^n \omega_i.
\]
The estimator is therefore translation invariant if and only if
\[
    \sum_{i=1}^n \omega_i = 0.
\]
In that case, adding a constant to all outcomes leaves the estimate unchanged:
\[
    \hat{\tau}(Y+c) = \hat{\tau}(Y).
\]

This condition explains the difference between the normalized and unnormalized
kappa estimators. Consider first a normalized estimator that takes the difference
between two normalized weighted means:
\[
    \hat{\tau}_{norm}
    =
    \frac{\sum_{i=1}^n a_i Y_i}{\sum_{i=1}^n a_i}
    -
    \frac{\sum_{i=1}^n b_i Y_i}{\sum_{i=1}^n b_i}.
\]
If a constant \(c\) is added to all outcomes, then
\[
    \hat{\tau}_{norm}(Y+c)
    =
    \frac{\sum_{i=1}^n a_i (Y_i+c)}{\sum_{i=1}^n a_i}
    -
    \frac{\sum_{i=1}^n b_i (Y_i+c)}{\sum_{i=1}^n b_i}.
\]
Expanding both terms yields
\[
    \hat{\tau}_{norm}(Y+c)
    =
    \left(
    \frac{\sum_{i=1}^n a_i Y_i}{\sum_{i=1}^n a_i}
    +
    c
    \right)
    -
    \left(
    \frac{\sum_{i=1}^n b_i Y_i}{\sum_{i=1}^n b_i}
    +
    c
    \right).
\]
The two constants cancel:
\[
    \hat{\tau}_{norm}(Y+c)
    =
    \frac{\sum_{i=1}^n a_i Y_i}{\sum_{i=1}^n a_i}
    -
    \frac{\sum_{i=1}^n b_i Y_i}{\sum_{i=1}^n b_i}
    =
    \hat{\tau}_{norm}(Y).
\]
Therefore, normalized estimators are translation invariant.

This applies to the normalized estimators reported in the replication, such as
\[
    \hat{\tau}_u
    =
    \frac{
    \frac{\sum_{i=1}^n Y_i Z_i / \hat{p}_i}
         {\sum_{i=1}^n Z_i / \hat{p}_i}
    -
    \frac{\sum_{i=1}^n Y_i (1-Z_i) / (1-\hat{p}_i)}
         {\sum_{i=1}^n (1-Z_i) / (1-\hat{p}_i)}
    }{
    \frac{\sum_{i=1}^n D_i Z_i / \hat{p}_i}
         {\sum_{i=1}^n Z_i / \hat{p}_i}
    -
    \frac{\sum_{i=1}^n D_i (1-Z_i) / (1-\hat{p}_i)}
         {\sum_{i=1}^n (1-Z_i) / (1-\hat{p}_i)}
    },
\]
where \(\hat{p}_i = \hat{P}(Z_i=1 \mid X_i)\). Since the outcome part consists of
a difference between normalized weighted means, adding \(\log(100)\) to all
log wages cancels from the numerator. Hence,
\[
    \hat{\tau}_u(Y+\log(100)) = \hat{\tau}_u(Y).
\]

The same logic applies to the normalized Abadie--Cattaneo estimator:
\[
    \hat{\tau}_{a,10}
    =
    \frac{\sum_{i=1}^n \hat{\kappa}_{1i}Y_i}
         {\sum_{i=1}^n \hat{\kappa}_{1i}}
    -
    \frac{\sum_{i=1}^n \hat{\kappa}_{0i}Y_i}
         {\sum_{i=1}^n \hat{\kappa}_{0i}}.
\]
After replacing \(Y_i\) by \(Y_i+c\), the constant appears once in each
normalized weighted mean and therefore cancels:
\[
    \hat{\tau}_{a,10}(Y+c)
    =
    \hat{\tau}_{a,10}(Y).
\]

The unnormalized estimators behave differently. For example, the unnormalized
Abadie estimator can be written as
\[
    \hat{\tau}_a
    =
    \frac{
    n^{-1}\sum_{i=1}^n
    Y_i
    \frac{Z_i-\hat{p}_i}{\hat{p}_i(1-\hat{p}_i)}
    }{
    n^{-1}\sum_{i=1}^n \hat{\kappa}_i
    }.
\]
If \(Y_i\) is replaced by \(Y_i+c\), then
\[
    \hat{\tau}_a(Y+c)
    =
    \frac{
    n^{-1}\sum_{i=1}^n
    (Y_i+c)
    \frac{Z_i-\hat{p}_i}{\hat{p}_i(1-\hat{p}_i)}
    }{
    n^{-1}\sum_{i=1}^n \hat{\kappa}_i
    }.
\]
Expanding the numerator gives
\[
    \hat{\tau}_a(Y+c)
    =
    \hat{\tau}_a(Y)
    +
    c
    \frac{
    n^{-1}\sum_{i=1}^n
    \frac{Z_i-\hat{p}_i}{\hat{p}_i(1-\hat{p}_i)}
    }{
    n^{-1}\sum_{i=1}^n \hat{\kappa}_i
    }.
\]
The second term is not generally equal to zero in finite samples. Therefore,
\[
    \hat{\tau}_a(Y+c)
    \neq
    \hat{\tau}_a(Y)
\]
in general. This means that the unnormalized estimator can change simply because
wages are recorded in cents rather than dollars before taking logarithms.

The same issue applies to the other unnormalized estimators reported in the
table, such as
\[
    \hat{\tau}_{a,1}
    =
    \frac{
    n^{-1}\sum_{i=1}^n
    Y_i
    \frac{Z_i-\hat{p}_i}{\hat{p}_i(1-\hat{p}_i)}
    }{
    n^{-1}\sum_{i=1}^n \hat{\kappa}_{1i}
    },
\]
and
\[
    \hat{\tau}_{a,0}
    =
    \frac{
    n^{-1}\sum_{i=1}^n
    Y_i
    \frac{Z_i-\hat{p}_i}{\hat{p}_i(1-\hat{p}_i)}
    }{
    n^{-1}\sum_{i=1}^n \hat{\kappa}_{0i}
    }.
\]
These estimators also do not generally have outcome weights that sum to zero.
Consequently, they are not generally translation invariant.

In the Vietnam draft lottery application, this distinction has a direct
empirical implication. The normalized estimators produce the same estimate
whether the log wage is constructed from dollars or cents:
\[
    \hat{\tau}_{norm}(\log w_i)
    =
    \hat{\tau}_{norm}(\log(100w_i)).
\]
By contrast, the unnormalized estimators can change:
\[
    \hat{\tau}_{unnorm}(\log w_i)
    \neq
    \hat{\tau}_{unnorm}(\log(100w_i)).
\]
Since dollars and cents describe the same underlying wage, this difference is
not economically meaningful. It is a finite-sample artifact of the estimator.
This is why normalization is important: it ensures that the estimated causal
effect is not affected by an arbitrary scaling choice in the outcome variable.

\end{document}
